\labsection{2}{Прогнозирование значений процесса с помощью линейной нейронной сети}

\labpart{Цель работы}
Изучить линейную нейронную сеть и метод наименьших квадратов для настройки её весов; построить
линейную сеть, прогнозирующую следующее значение дискретного сигнала по $k$ предыдущим значениям, и
исследовать влияние шага дискретизации и размера окна на точность прогноза.

\labpart{Теоретические сведения}
Подразделы~\ref{sec:perceptron}.4 и~\ref{sec:perceptron}.5 теоретической части: линейная сеть ADALINE,
решение через псевдообратную матрицу (\ref{eq:pinv}), линия задержки и матрица входов
(\ref{eq:delay-matrix}), рекуррентное соотношение для гармонического сигнала (\ref{eq:harmonic}),
теорема Котельникова.

\labpart{Постановка задачи}
Сигнал $x(t) = f(t)$ задан на отрезке $t \in [T_1; T_2]$ и дискретизирован с шагом $0 < \Delta t \le 0,05$~с.
Необходимо построить линейную нейронную сеть, прогнозирующую будущее значение сигнала по $3 \le k \le 7$
предыдущим значениям, получить вектор прогнозов $y$ моделированием сети и построить:
график сигнала $y(t)$, на котором символом <<\code{*}>> отмечены точки эталонного сигнала $x(t)$, и
график ошибки $e = y - x$. Варьируя $\Delta t$ и $k$, минимизировать ошибку. Варианты заданий приведены
в таблице~\ref{tab:l2-var}: для синусоидальных сигналов $t \in [0; 5]$, для косинусоидальных
$t \in [-3; 3]$.

\begin{table}[!htb]
\caption{Варианты заданий к лабораторной работе №2}
\label{tab:l2-var}
\centering\small
\begin{tabular}{cc|cc|cc|cc}
\toprule
№ & $x(t)$ & № & $x(t)$ & № & $x(t)$ & № & $x(t)$ \\
\midrule
1 & $\sin(\pi t)$     & 5 & $\sin(2\pi t)$   & 9  & $\sin(3\pi t)$   & 13 & $\sin(4\pi t)$ \\
2 & $\cos(\pi t)$     & 6 & $\cos(2\pi t)$   & 10 & $\cos(3\pi t)$   & 14 & $\cos(4\pi t)$ \\
3 & $\sin(1,5\pi t)$  & 7 & $\sin(2,5\pi t)$ & 11 & $\sin(3,5\pi t)$ &    & \\
4 & $\cos(1,5\pi t)$  & 8 & $\cos(2,5\pi t)$ & 12 & $\cos(3,5\pi t)$ &    & \\
\bottomrule
\end{tabular}
\end{table}

\labpart{Порядок выполнения работы}
\begin{steps}
\item Для заданного отрезка и шага $\Delta t$ сформировать вектор времени $t$ и вектор значений сигнала $x$.
\item Сформировать матрицу входов $P$ размера $k \times N$ (\ref{eq:delay-matrix}): $i$-й столбец содержит
$k$ значений сигнала, предшествующих $x_i$, начальные значения линии задержки нулевые. Целевой
вектор~--- $x$.
\item Найти веса линейной сети с $k$ входами и одним выходом по МНК (\ref{eq:pinv}):
\code{np.linalg.lstsq} или \code{np.linalg.pinv} с добавлением строки единиц для смещения
(подраздел~\ref{sec:py-perceptron}); для контроля сравнить с \code{sklearn.linear\_model.LinearRegression}.
\item Вычислить прогноз сети $y = WP + b$.
\item Построить графики $x(t)$, $y(t)$ и ошибки $e(t)$ в отдельных окнах; вычислить MSE и максимальную
ошибку.
\item Исследовать зависимость ошибки от $k \in [3; 7]$ и $\Delta t \in (0; 0,05]$; сравнить найденные
веса с теоретическими значениями из соотношения (\ref{eq:harmonic}).
\item Дополнительно: добавить к сигналу нормальный шум ($\sigma = 0,05$) и повторить исследование
зависимости ошибки от $k$, оценивая ошибку относительно незашумлённого сигнала.
\end{steps}

\labpart{Пример выполнения}
Вариант~3: $x(t) = \sin(1,5\pi t)$, $t \in [0; 5]$, $\Delta t = 0,05$~с ($N = 101$ отсчёт), $k = 5$. МНК
дал веса $\mathbf{w} = (1,934;\ -0,990;\ -0,001;\ 0,010;\ -0,011)$ и смещение $b = 0,002$. Первые два веса
близки к теоретическим коэффициентам $2\cos(\omega\Delta t) = 1,945$ и $-1$ из соотношения
(\ref{eq:harmonic}), остальные~--- к нулю. Прогноз практически совпадает с сигналом
(рисунок~\ref{fig:l2-pred}), MSE $= 5,3 \cdot 10^{-4}$. Почти вся ошибка сосредоточена в начале ряда,
пока линия задержки заполнена нулями: отсчёт $x(\Delta t) = 0,233$ не может быть предсказан по нулевой
предыстории.

\begin{figure}[!htb]
\centering
\includegraphics[width=\textwidth]{\figdir/lab2_prediction.png}
\caption{Прогноз линейной сети для варианта~3 ($\Delta t = 0,05$, $k = 5$) и ошибка прогноза}
\label{fig:l2-pred}
\end{figure}

Исследование параметров (рисунок~\ref{fig:l2-study}) показывает, что для незашумлённой синусоиды
ошибка не зависит от $k$ при $k \ge 2$ и совпадает с оценкой $\sin^2(\omega\Delta t)/N$: уменьшение шага
с $0,05$ до $0,005$ снижает MSE в тысячу раз. Для зашумлённого сигнала картина иная: при $k = 2$ коэффициенты
$2\cos(\omega\Delta t) \approx 2$ и $-1$ усиливают шум, и ошибка до четырёх раз превышает его дисперсию
$\sigma^2 = 0,0025$. С ростом $k$ прогноз усредняет шум по большему числу отсчётов, и ошибка
уменьшается в 2,5--5 раз; при $\Delta t = 0,05$ насыщение наступает около $k = 6$--$7$.

\begin{figure}[!htb]
\centering
\includegraphics[width=\textwidth]{\figdir/lab2_study.png}
\caption{Зависимость ошибки прогноза от шага дискретизации (чистый сигнал) и от размера окна (сигнал с шумом)}
\label{fig:l2-study}
\end{figure}

\labpart{Контрольные вопросы}
\begin{questions}
\item Чем линейная сеть отличается от персептрона? Какую функцию ошибки она минимизирует?
\item Запишите правило Уидроу~--- Хоффа. Чем ограничен шаг обучения?
\item Выведите формулу весов линейной сети через псевдообратную матрицу. Почему \code{np.linalg.lstsq}
предпочтительнее явного обращения матрицы $PP^{\mathsf T}$?
\item Как формируется матрица входов для прогнозирования временного ряда? Что такое линия задержки?
\item Почему для прогноза синусоиды достаточно двух предыдущих значений? Выведите соотношение
(\ref{eq:harmonic}).
\item Где на графике ошибки сосредоточена основная ошибка прогноза и почему?
\item Как шаг дискретизации влияет на точность прогноза? Сформулируйте теорему Котельникова.
\item Как влияет размер окна $k$ на прогноз зашумлённого сигнала?
\item Какие сигналы невозможно точно прогнозировать линейной сетью?
\end{questions}
